\chapter{Análisis de estudios observacionales y RCTs}
Vamos a aplicar el marco presentado primero para estudios randomizados controlados y
posteriormente para estudios observacionales.
Vamos a ver las condiciones que un estudio randomizado implica y cómo se debilitan
en un estudio observacional.

\section{Estudios randomizados}
Vamos a definir unos conceptos más fuertes que la intercambiabilidad condicional:

\begin{definicion}[Randomización]
    Diremos que un estudio implicando variables $X, A, Y$ es randomizado cuando la
    variable de tratamiento $A$ es controlada. En el caso más sencillo, exigimos independencia,
    $P(A | X) = P(A)$.
    % $A \indep Y(a)$: el tratamiento no está correlado con otras covariantes ni % TODO:
\end{definicion}

\begin{nota}
    Sólo consideramos el caso en el que $P(A) = P(A | X)$, con un muestreo simple sin estratificación
    ni otras medidas más sofisticadas.
\end{nota}

\begin{definicion}[Intercambiabilidad]
    Decimos que se presenta intercambiabilidad cuando $ Y(a) \indep A, \, \forall a $.
    Escrito de otra forma, $P(Y(a) = y | A = 1) = P(Y(a) = y | A = 0) = P(Y(a) = y)$,
    en el caso discreto.
\end{definicion}

\begin{proposicion}[Randomización implica intercambiabilidad]

\end{proposicion}

\begin{proposicion}[Intercambiabilidad implica intercambiabilidad condicional]

\end{proposicion}

Consideramos de nuevo el grafo causal de un estudio randomizado controlado:

\begin{figure}[h]
    \label{fig:grafo-causal-rct}
    \centering
    \begin{tikzcd}
        A \arrow{r}{} & Y \\
        X \arrow{ur}{}
    \end{tikzcd}
    \caption{RCT}
\end{figure}

Habiendo varias formas de randomizar a los sujetos\footnotemark, consideramos una
randomización simple en la que los sujetos son asignados a las variantes del tratamiento
por una distribución fija sin parámetros libres.

% TODO:
\begin{proposicion}[Propiedades de un estudio controlado]
    Se verifican las siguientes propiedades:
    \begin{itemize}
        \item $A \indep X$
        \item  $A \indep Y(a)$
    \end{itemize}
\end{proposicion}


\footnotetext{Una randomización simple como la usada en nuestro ejemplo puede no ser robusta
    frente a muestras pequeñas. Existen técnicas como la randomización por bloques,
    que divide la población a estudiar en estratos
    por alguna variable que no sea de interés primario (por ejemplo, el sexo o la edad)
    y asignando el tratamiento de forma randomizada en cada estrato.
    De esta forma, uno controla la variabilidad que hubiera generado tal variable.
    A estas variables se las denomina, en inglés, \say{nuisance variables}. En el grafo causal,
    $X$ es una de ellas.
}

\section{Estudios observacionales}
Ahora, no controlamos la asignación del tratamiento, sino que puede ser afectada por el resto de
condiciones del sujeto de estudio ($X$).

\begin{figure}[h]
    \label{fig:grafo-causal-obs}
    \centering
    \begin{tikzcd}
        A \arrow{r}{} & Y \\
        X  \arrow{u}{} \arrow{ur}{}
    \end{tikzcd}
    \caption{Estudio observacional}
\end{figure}
